% --- KOPIEER DIT BESTAND VOOR ELKE NIEUWE DAG ---

\subsection*{Wanneer \& wie?: [09-06-2026] - [Jonas]}

\textbf{Wat heb je gedaan?}\\
Errorbars toegevoegd op onze berekende magnitude, deze errorbars hebben we berekend  met de poisson formule, hoewel ik niet de onzekerheid op de magnitude van szlyn heb berekend, heb ik dit gedaan op de (ref-check) magnitude, aangezien dit overeenkomt. 

Ook heb ik gekeken naar het schalen van de dark-frames. Gezien wij op een nacht alleen maar dark frames van 50s hebben en niet van 5s (de exptime van flats) moeten we ze schalen.
\vspace{0.3cm} 

\noindent \textbf{Waarom heb je dat gedaan?}\\
Aangezien wij op onze fase-diagram een systematische jump zien tussen dagen van observeringen, dit zou een continue lijn moeten zijn.

Wellicht kwam dit door het niet correct opstellen van de master-flats door die exptime.

\noindent \textbf{Dingen die noemenswaardig zijn:}\\
Het is me gelukt om errorbars te krijgen maar die vormen geen duidelijkheid over de jump, ook heb ik gekeken naar de masterflats die ik nu volgensmij goed heb geschaald, heb het als een opmerking voor de meeting van morgen neergezet. Zodat ik sebastian hierover kan vragen.


\vspace{0.5cm}
\hrule 
\vspace{0.5cm}
